% Source: https://artofproblemsolving.com/community/c3h1608289_2010_national_mathcounts_questions_urgent
% Stable identity: aops-topic-1608289-q05
\textbf{The boxes count in base $5$.} Read the numbers of balls in the boxes, from right to left, as the digits of a number; each digit is between $0$ and $4$. A step adds one ball to the first box that is not full and empties the full boxes before it. That is exactly how $1$ is added to a number in base $5$: trailing digits $4$ become $0$ and the first digit less than $4$ increases by $1$. So after step $n$ the boxes show $n$ written in base $5$ (the first box holding the units digit).

\textbf{Write $2010$ in base $5$.} $2010=3\cdot625+135$, $135=1\cdot125+10$, $10=0\cdot25+2\cdot5+0$, so
\[
2010=3\cdot5^4+1\cdot5^3+0\cdot5^2+2\cdot5+0=31020_5 .
\]
The total number of balls is the digit sum $3+1+0+2+0=\boxed{6}$.
